\chapter{机器学习研究设计：验证、监控与 Alpha 决策}
\label{ch:lower-ml-validation}

\begin{itemize}
  % Generated from curriculum/manifest.json; do not edit by hand.
\item 直接先修：下册《表格统计学习：从决策树到梯度提升》。

\end{itemize}

排序正确与概率正确是不同目标。一个模型可以把所有上涨样本排在前面，却把 60\% 的上涨概率报成 99\%；下面把时间验证与概率评价分别计算。


\section{嵌套时间轴与选择权}

按信息完成时间建立训练与验证，而不只按行号。若第 $t$ 行标签在 $t+h$ 才完成，拟合时只能使用已完成标签。以下为一种保守的分块选择：
\[
 \max I_{\rm train}+h<\min I_{\rm val},\qquad
 \max I_{\rm val}+h<\min I_{\rm test}.
\]
若另外需要缓冲 $e$，应在标签完成时间之外再加 $e$。例如训练最后索引 7，$h=2$，验证从 10 开始；验证最后索引 11，其标签到 13 完成，最终测试从 14 开始。额外缓冲并非通用必需条件，不能保证数据独立。

% Generated by tools/build_figures.py; edit figures/figure-specs.json instead.
\MFQFigure{tex/figures/generated/lower-ml-nested-timeline.pdf}{内层窗口用于模型选择，外层窗口评价按这一选择得到的预测；汇总外层结果即可比较整套学习方法。}{lower-ml-nested-timeline}


内层只能在外层训练区内选择超参数；外层测试衡量整个选择程序。合法特征只可由决策时点的信息集 $\mathcal I_t$ 构造。把测试结果反馈给模型后，它就变成新的验证集。尝试次数、搜索空间和提前停止规则都属于研究选择权计算记录。

\section{Cross-fitting 与概率校准}

cross-fitting 对每个验证块只使用更早训练块拟合，再把预测写回原时间索引。透明 NumPy 路径显式构造截距、只惩罚斜率并拒绝倒序或重叠折；成熟库路径使用相同折和相同 ridge 强度的 \texttt{Ridge}。二者事先确定预测差低于数值容差。前者适合手算与暴露约束，后者适合扩展特征维度；若库默认截距、标准化或惩罚约定不同，就不能把差异归咎于浮点误差。cross-fitting 减少同一样本同时拟合 nuisance model 和评价主效应的污染，但不能自动解决非平稳或标签泄漏。事先确定两折产生索引 4--7 的四个真正样本外预测。

% Generated by tools/build_figures.py; edit figures/figure-specs.json instead.
\MFQFigure{tex/figures/generated/lower-ml-calibration.pdf}{校准图将预测概率与样本中的事件频率比较；排序能力相近的模型，概率数值仍可能有明显差异。}{lower-ml-calibration}


概率校准拟合 $\Pr(Y=1\mid s)$，例如 Platt scaling 的 logistic 映射。Brier 分数
\[
\operatorname{BS}=n^{-1}\sum_i(p_i-y_i)^2
\]
在较早校准窗口拟合 logistic 映射，再在后期评价窗口计算分数；绝不在同一批标签上自评改善。事先确定后期 Brier 的实际数值由 notebook 重建。校准改善仍不代表未来可靠；可靠性图、样本数与置信区间应同时报告\cite{niculescu2005probabilities}。

\MFQTransition{偏差、方差与金融样本的有效规模}

监督学习把泛化误差分解为不可约噪声、模型偏差和估计方差只是起点。金融面板中的行数 $N\times T$ 不能直接当作独立样本数：同一时点资产共享市场冲击，同一资产跨期相关，重叠标签又共享未来收益。模型容量应相对于独立时间块和横截面簇来判断。

学习曲线要同时沿两个方向画：增加历史长度和增加横截面资产。若只增加同一短时期的资产就改善，模型可能主要学习截面结构；若跨年份扩展后性能下降，说明制度漂移比方差下降更重要。训练误差很低、时序验证误差高是过拟合；两者都高则可能是特征、目标或容量不足。

树深、boosting 轮数、网络宽度、dropout、学习率和 lookback 都是选择权。包含 20 个候选配置、5 个随机种子和 4 种标签的研究实际上做了 400 次尝试。报告只展示获胜模型会隐藏 data snooping；至少保留候选空间、失败结果和停止规则。

\MFQTransition{嵌套 walk-forward 的完整算法}

设外层评价窗口为 $O_k$，其之前的开发数据为 $D_k$。在 $D_k$ 内再建立若干按时间前推的内层训练/验证折，选择超参数 $\lambda_k$；随后用整个 $D_k$ 重新拟合并只预测 $O_k$。不同外层折允许因市场变化选择不同超参数，但最终报告必须聚合整个选择程序，而不是从每折候选中事后再选一次。

标签 $y_t=r_{t+1:t+h}$ 跨越 $h$ 期时，任何训练样本的标签区间与验证区间相交都应 purge。Embargo 是验证/测试边界后的额外空窗，用于阻断特征或标签的缓慢传播。它们按信息区间定义，不是机械删除固定行数。

预处理器也在折内拟合：缺失填充、winsorize、标准化、PCA、目标编码、特征筛选都只能读取训练子集。把全样本 scaler 放在 pipeline 外面，是最常见且最隐蔽的泄漏之一。

\section{指标、解释与漂移响应}

回归报告 MSE、MAE、$R^2_{\mathrm{OOS}}$；横截面排序报告 Pearson IC、Rank IC、Top--Bottom spread 和分组单调性；分类报告 ROC-AUC、PR-AUC、log loss 与 Brier。类别极不平衡时 ROC-AUC 可能看似良好而精确率很差。

% Generated by tools/build_figures.py; edit figures/figure-specs.json instead.
\MFQFigure{tex/figures/generated/lower-ml-drift.pdf}{输入分布、预测分布和条件收益关系可能分别变化；区分这些变化有助于判断重新估计哪些部分。}{lower-ml-drift}


经济指标必须建立在事先确定的仓位映射后，包括净收益、波动、回撤、换手、成本、容量和暴露。模型 A 的 MSE 较低不保证组合收益更高：改善可能集中在接近零的样本，而排序头尾没有变化。评价单位应匹配决策；逐行置信区间会忽略同日相关，通常要按日期聚合或使用双向聚类、块 bootstrap。

\MFQTransition{校准、阈值与拒绝决策}

概率模型的 $0.7$ 应意味着在相似条件下约 70\% 发生，而不只是比 $0.6$ 排名更高。Platt scaling 假设 logistic 映射，isotonic 更灵活但小样本易过拟合。校准器只能在模型未用于拟合的预测上训练，常用 cross-fitted scores。

若只有预测概率超过 $q$ 才交易，$q$ 必须在验证期结合成本和容量选择。更稳健的系统允许拒绝：当不确定性过大、输入漂移或关键特征缺失时输出“无交易”，而不是强制每个样本产生方向。

\MFQTransition{解释：全局结构、局部归因与反事实}

树的 split gain、permutation importance、PDP/ALE、SHAP 和神经网络梯度回答不同问题。相关特征下，单列置换会制造不现实样本；PDP 也可能穿越数据支持之外。ALE 在局部条件区间累积效应，通常更适合相关特征，但仍是模型解释而非因果效应。

局部归因的加和恒等式只能证明解释器分配了预测差，不证明分配唯一或经济正确。应改变训练窗口、随机种子和背景样本检查重要性稳定性，并用删除特征组后的样本外损失做消融。财务特征存在会计恒等式，把市值单独减半却保持股价和股本不变不是合法反事实。

\MFQTransition{解释稳定性与漂移响应}

Top-$k$ 重要特征集合的 Jaccard 衡量两个窗口使用的特征是否一致。本章 Top-2 只有 $1/3$，必须报告不稳定；即使为 1，也不证明因果。正式研究可加入 permutation importance、SHAP 或消融，但每种方法都要声明背景分布和相关特征处理。

漂移至少区分输入漂移 $P(X)$、标签漂移 $P(Y)$、概念漂移 $P(Y\mid X)$ 与执行漂移。事先确定均值距离为 $0.75$，超过 0.5 后停止自动上线并进入重校准/重训练审批。响应动作需事先确定，不能看到亏损后再挑指标解释。

输入漂移可用 PSI、KS、Wasserstein 距离或分类器区分新旧样本；标签漂移要等标签成熟后观察；概念漂移比较条件误差或残差结构；执行漂移来自成交、成本和覆盖率变化。多指标持续监控本身又形成多重检验，阈值和联合报警规则需事前设定。

报警后的动作可分级：记录但不行动、降低仓位、回退基线、重新校准、重新训练或停止模型。重新训练不能自动吸收刚发生的异常期；需要最小样本、审批与新的外层留出。监控面板应同时显示数据质量、模型输出、组合暴露和执行状态，避免把所有问题都归咎于模型。

\begin{heuristic}
漂移统计量回答“分布是否变化”，不是“变化是否有经济意义”。极显著但幅度很小的漂移可能无关紧要；样本少但影响成交规则的制度变化可能必须立即停用。
\end{heuristic}

本单元向路线 Capstone 交付：嵌套切分、cross-fitted 预测、校准报告、解释稳定性和漂移响应记录。
\subsection{Brier 分数为何鼓励诚实概率}

设真实成功概率为 $q$，预测为 $p$。条件期望损失为
\[
 \mathbb E(p-Y)^2=q(p-1)^2+(1-q)p^2
 =(p-q)^2+q(1-q).
\]
第二项与预测无关，第一项在且仅在 $p=q$ 时最小，所以 Brier 是严格适当评分。若 $q=0.6$，报告 $p=0.6$ 的期望损失为 $0.24$，报告 $0.9$ 为 $0.33$；更极端不等于更好。

有限评价样本 $p=(0.2,0.8,0.6,0.4)$、$y=(0,1,1,0)$ 的 Brier 为 $(0.04+0.04+0.16+0.16)/4=0.1$。把概率分箱后比较频率可以展示校准，但每箱样本少时波动很大。notebook 用较早数据拟合预测，后续数据拟合校准器，再在更晚数据评价；不能在同一组标签上调概率又宣称校准改善。
